% Richard Lin – Resume LaTeX Template
% Compile with: lualatex richard_lin_resume.tex

\documentclass[11pt, letterpaper]{article}

% ── Packages ────────────────────────────────────────────────────────────────
\usepackage[margin=0.5in, top=0.5in, bottom=0.35in]{geometry}
\usepackage{enumitem}
\usepackage{titlesec}
\usepackage{hyperref}
\usepackage[sfdefault]{carlito}
\setsansfont[
    Ligatures      = {NoCommon, NoRequired, NoContextual},
    UprightFont    = *-Regular,
    ItalicFont     = *-Italic,
    BoldFont       = *-Bold,
    BoldItalicFont = *-BoldItalic
]{Carlito}
\renewcommand{\familydefault}{\sfdefault}
\usepackage{microtype}

\hypersetup{
    colorlinks=true,
    urlcolor=black,
    linkcolor=black
}

\titleformat{\section}
    {\normalfont\fontsize{12}{14}\selectfont\bfseries\uppercase}
    {}{0em}{}
    [\titlerule]
\titlespacing{\section}{0pt}{4pt}{2pt}

\setlist[itemize]{leftmargin=1.5em, topsep=0pt, itemsep=0pt, parsep=0pt}

\pagestyle{empty}

\setlength{\parskip}{0pt}
\setlength{\parsep}{0pt}

\newcommand{\roleentry}[4]{%
    \noindent\textbf{#1} \hfill #2 \\
    \textit{#3} \hfill \textit{#4}%
    \vspace{0pt}
}

\newcommand{\subrole}[2]{%
    \noindent\textit{#1} \hfill \textit{#2}%
    \vspace{0pt}
}

\newcommand{\projectentry}[2]{%
    \noindent\textbf{#1} \hfill \textit{#2}%
    \vspace{0pt}
}

\begin{document}

\noindent
\hbox to \textwidth{\hfil{\fontsize{16}{19}\selectfont\textbf{Richard Lin}}\hfil\makebox[0pt][r]{U.S. Citizen}}\\[-1pt]
{\centering
    Huntingdon Valley, PA | 215-403-4999 |
    \href{mailto:richard.rich.lin@gmail.com}{richard.rich.lin@gmail.com} \\
    \href{https://github.com/richardlin05}{github.com/richardlin05} |
    \href{https://richardlin05.github.io}{richardlin05.github.io} |
    \href{https://linkedin.com/in/richardlin05}{linkedin.com/in/richardlin05}
    \par}

\vspace{-2pt}

\section{Education}

\roleentry{Temple University}{Philadelphia, PA}%
    {Bachelor of Science in Electrical and Computer Engineering}{Expected Graduation: May 2028}

\noindent Minor: Computer Science | Fundamentals of Physics Certificate | GPA: 3.88/4.00 | Dean's List: F24, S25, F25 \\
\noindent \textbf{Courses:} Classical Control Systems; Signals \& Systems; Processor Systems; Engineering Computation IV (GPU Programming)

\vspace{2pt}

\section{Professional Experience}

\roleentry{Temple University}{Philadelphia, PA}%
    {Undergraduate Research Assistant -- Computer Fusion Lab (CFL)}{Oct 2025 -- Present}

\noindent \textit{Advisor: Dr. Li Bai}
\begin{itemize}
    \item Administered lab website and systems; ensured security, software maintenance and network reliability
    \item Built a multi-agent LLM system in Python: a service agent gathers a patron's situation, then a library agent (RAG + LLM-rerank) recommends top 3 best-fit devices from a 400+ item catalog (synced daily), helping people with disabilities
    \item Trained a Vision-Language-Action (VLA) model with LeRobot to enable a robotic arm to pick up and hand objects for individuals with limited mobility; developed ROS-based robot-control in C++ across multiple robots
\end{itemize}

\vspace{1pt}
\subrole{Undergraduate Teaching Assistant}{Aug 2025 -- Present}
\begin{itemize}
    \item Assisted faculty with instruction, hands-on activities, and assignment evaluation; held office hours for student support
    \item \textbf{Intro to Engineering:} Guided students building embedded Micro:bit incubator, testing components (LEDs, fans, sensors)
    \item \textbf{Engineering Problem Solving:} Supported wind turbine calculation/data analysis using Excel, Python (NumPy, Matplotlib)
    \item \textbf{Engineering Computation I:} Taught C/C++/Python/Linux/Bash fundamentals, plus basic algorithms and time complexity
\end{itemize}

\vspace{1pt}
\subrole{STEPS Ambassador -- Sustainable Temple Energy and Power Scholars (STEPS)}{Sep 2024 -- Present}
\begin{itemize}
    \item Mentored STEPS underclassmen through the high-school-to-college transition, providing academic guidance
    \item Designed and led renewable-energy STEM demos for high schoolers, emphasizing sustainability and real-world engineering
\end{itemize}

\vspace{2pt}

\section{Projects}

\projectentry{Musical Light Show -- ECE 2352: Circuits and Electronics II}{May 2026}
\begin{itemize}
    \item Designed and built an analog audio-processing system that combines left/right audio channels, and separates high and low frequency signals to drive an audio-reactive LED light show
    \item Implemented dual op-amp stages for summing and volume/gain control using a feedback potentiometer, high/low-pass filters, peak detectors, BJT LED driver circuits
    \item Calculated values from gains and cutoff frequency using Ohm's Law and KVL; verified in Multisim and hardware within 1\%
\end{itemize}

\vspace{1pt}
\projectentry{Detective Conan AI Glasses (Childhood Dream Project)}{Aug 2026 -- Present}
\begin{itemize}
    \item Architecting smart glasses on Pi Zero 2 W with camera, mic, solar charging, 3D-printed enclosure; Google Gemini API
    \item Developing software architecture with intent, LLM routing, object detection, OCR, face recognition, and cybersecurity
\end{itemize}

\vspace{2pt}

\section{Clubs \& Leadership}

\roleentry{Institute of Electrical and Electronics Engineers (IEEE) -- Temple University}{Philadelphia, PA}%
    {President}{Sep 2024 -- Present}
\begin{itemize}
    \item Led 11-member IEEE executive board in planning outreach programs, technical workshops, and industry networking events, while coordinating with four affiliated chapters and societies (EMBS, PES, WIE, IEEE-HKN)
    \item Grew IEEE SAC participation by 50\% (14 → 21) through targeted event promotion and outreach during 2025–2026
    \item Founded the first annual multi-university hardware hackathon with Drexel and UPenn IEEE branches
\end{itemize}

\vspace{1pt}
\roleentry{Eta Kappa Nu (IEEE-HKN) -- Temple University}{Philadelphia, PA}%
    {Secretary \& Treasurer}{Jun 2026 -- Present}

\vspace{1pt}
\roleentry{Temple University College of Engineering Alumni Association}{Philadelphia, PA}%
    {Student Representative}{Jun 2026 -- Present}

\vspace{2pt}

\section{Skills}

\noindent
\textbf{Programming:} Python, C, C++, Java, MATLAB, Bash, SystemVerilog, JavaScript, HTML/CSS, CUDA, AVR Assembly \\[1pt]
\textbf{AI \& Robotics:} Machine Learning, PyTorch, LLM Integration, RAG, VLA, Edge AI, Computer Vision (OpenCV, YOLO), ROS \\[1pt]
\textbf{Hardware \& Tools:} Git/GitHub, Linux, Raspberry Pi, Arduino, FPGA, Circuit Design \& Analysis, Soldering, n8n, PCB Design \\[1pt]
\textbf{Languages:} English, Mandarin Chinese, Fuzhounese, Spanish

\end{document}